\section{学习型大气初始化 Learned Atmosphere Initialization}\label{sec:atmosphere_emulator}

我们通过学习初始状态来削减上述剩余成本，同时保留物理修正循环。一趟循环重建压强与布居、不透明度、辐射转移、对流以及重映射后的温度和柱质量结构。它不构建任何全局导数矩阵，因此每次迭代的成本就是一趟物理遍历。这种经济的更新方式在热矮星、太阳、红巨星和 K 矮星测试中表现良好。在我们的测试中，更通用的求根算法虽然提高了稳健性，但需要多得多的物理遍历。

\begin{EN}
We reduce this remaining cost by learning the starting state while retaining
the physical correction cycle. One cycle
rebuilds pressure and populations, opacity, transfer, convection, and the
remapped temperature and column-mass structure. It forms no global derivative
matrix, so each iteration costs one physical pass. This economical update
works well across the hot-dwarf, solar, red-giant, and K-dwarf tests. In our
tests, more general root solvers improved robustness but required many physical
passes.
\end{EN}

收敛与否仍取决于从初始状态出发的反复更新是否趋近物理解。因此我们用\textbf{物理重启}（而非仅看廓线精度）来评判初始化器。所有比较都保持目标标签与丰度不变，只改变初始大气。预测结果不会被直接当作大气接受：物理求解器必须通过其收敛检查，另有独立检验考察流量闭合与合成光谱。

\begin{EN}
Convergence still depends on whether repeated updates from the starting state
approach the physical solution. We therefore judge an initializer by its
physical restart rather than profile accuracy alone. Every comparison keeps
the target labels and abundances fixed and changes only the starting atmosphere.
The prediction is not accepted as an atmosphere. The physical solver must pass
its convergence checks, while independent tests examine flux closure and the
synthesized spectrum.
\end{EN}

\begin{physbox}{学习型初始化器的定位：只猜初值，不做结论}
这是本文方法论中最“克制”的一处机器学习应用。神经网络不预测光谱、不替代物理，只做一件事：根据恒星标签猜一个足够好的大气初始结构，让物理迭代少跑几趟。最终结果被接受的充要条件仍是物理求解器收敛、流量平衡闭合、光谱检验通过。这与 The Payne 等“标签直接到光谱”的模拟器有本质区别：模拟器的误差会直接进入最终科学产物，而初始化器的误差只会影响收敛速度，不影响收敛到的物理解本身（只要初值落在收敛域内）。这体现了一条重要原则：\textbf{让机器学习加速探索，让物理负责裁决}。
\end{physbox}

\begin{notebox}{表~\ref{tab:initializer_families} 说明（原文 Table 4，内容保留英文）}
三个初始化器家族：\textbf{五标签}（$T_{\rm eff}$、$\log g$、$[\mathrm{M}/\mathrm{H}]$、$[\alpha/\mathrm{M}]$、$\xi$；52,199 个收敛大气）；\textbf{八标签 CNO}（增加 C、N、O 独立丰度；53,824 个）；\textbf{直接丰度}（$T_{\rm eff}$、$\log g$、$\xi$ 加 81 个独立 $[\mathrm{X}/\mathrm{H}]$，内部以 $[\mathrm{Fe}/\mathrm{H}]$ 加 80 个非铁 $[\mathrm{X}/\mathrm{Fe}]$ 表示；82,016 个）。$N_{\rm corpus}$ 为划分训练/验证/测试前保留的独立收敛大气数。
\end{notebox}

\begin{table*}[t]
\centering
\caption{Learned atmosphere initializer families, sample sizes, input labels,
and training envelopes. $N_{\rm corpus}$ is the number of independently
converged atmospheres retained before splitting. The next column gives the
training, validation, and held-out test counts. The public direct-abundance
input contains 81 $[\mathrm{X}/\mathrm{H}]$ values, represented internally by
$[\mathrm{Fe}/\mathrm{H}]$ and 80 non-iron $[\mathrm{X}/\mathrm{Fe}]$
coordinates.}
\label{tab:initializer_families}
\begingroup
\setlength{\tabcolsep}{5pt}
\renewcommand{\arraystretch}{1.20}
\footnotesize
\begin{tabular}{@{}lllll@{}}
\toprule
\tcell{0.11\textwidth}{\textbf{Family}} &
\tcell{0.08\textwidth}{$N_{\rm corpus}$} &
\tcell{0.16\textwidth}{\textbf{Training / validation / test}} &
\tcell{0.27\textwidth}{\textbf{Input labels}} &
\tcell{0.30\textwidth}{\textbf{Training envelope}} \\
\midrule
\tcell{0.11\textwidth}{\textbf{Five-label}} &
\tcell{0.08\textwidth}{52,199} &
\tcell{0.16\textwidth}{50,000 / 2,000 / 199} &
\tcell{0.27\textwidth}{$T_{\rm eff}$, $\log g$, $[\mathrm{M}/\mathrm{H}]$,
$[\alpha/\mathrm{M}]$, $\xi$} &
\tcell{0.30\textwidth}{$4000\leq T_{\rm eff}\leq10500$ K,
$0.7\leq\log g\leq5.3$, $-2.5\leq[\mathrm{M}/\mathrm{H}]\leq0.5$,
$-0.1\leq[\alpha/\mathrm{M}]\leq0.5$,
$0.5\leq\xi\leq4.0$ km s$^{-1}$} \\
\tcell{0.11\textwidth}{\textbf{Eight-label CNO}} &
\tcell{0.08\textwidth}{53,824} &
\tcell{0.16\textwidth}{51,625 / 2,000 / 199} &
\tcell{0.27\textwidth}{$T_{\rm eff}$, $\log g$, $[\mathrm{M}/\mathrm{H}]$,
$[\alpha/\mathrm{M}]$, $\xi$, $[\mathrm{C}/\mathrm{M}]$,
$[\mathrm{N}/\mathrm{M}]$, $[\mathrm{O}/\mathrm{M}]$} &
\tcell{0.30\textwidth}{Five-label bounds, with $-0.5\leq[\mathrm{C}/\mathrm{M}],
[\mathrm{N}/\mathrm{M}],[\mathrm{O}/\mathrm{M}]\leq0.5$} \\
\tcell{0.11\textwidth}{\textbf{Direct abundance}} &
\tcell{0.08\textwidth}{82,016} &
\tcell{0.16\textwidth}{71,660 / 9,733 / 623} &
\tcell{0.27\textwidth}{$T_{\rm eff}$, $\log g$, $\xi$, and 81
$[\mathrm{X}/\mathrm{H}]$ abundances, represented by
$[\mathrm{Fe}/\mathrm{H}]$ and 80 non-iron
$[\mathrm{X}/\mathrm{Fe}]$ coordinates} &
\tcell{0.30\textwidth}{$4000\leq T_{\rm eff}\leq10500$ K,
$0.7\leq\log g\leq5.3$, $0.5\leq\xi\leq4.0$ km s$^{-1}$,
$-2.5\leq[\mathrm{Fe}/\mathrm{H}]\leq0.5$, and
$-0.5\leq[\mathrm{X}/\mathrm{Fe}]\leq0.5$} \\
\addlinespace[2pt]
\bottomrule
\end{tabular}
\endgroup
\end{table*}

\begin{figure}[t]
\centering
\includegraphics[width=0.98\textwidth]{figure_06.pdf}
\caption{五标签与八标签初始化器的精度与迭代成本。左图比较 $K=160$ PCA 压缩下限（虚线、空心）与完整的留出标签到大气误差（实线、实心）；中图给出六个对照在三趟物理求解后的保守光谱误差上限（均低于 $5.1\times10^{-3}$）；右图比较“邻近存档大气”起点与两种学习型起点的迭代次数，横线标记中位数约 12、5.5、4 趟。（Accuracy and iteration cost of the five- and eight-label atmosphere initializers.）}
\label{fig:initializer_convergence}
\end{figure}

\subsection{作为不动点问题的初始化 Initialization as a fixed point problem}

固定所请求的 $T_{\rm eff}$、$\log g$、$\xi$ 与丰度 $\{[\mathrm{X}/\mathrm{H}]\}$，令 $\mathbf{x}^{(n)}$ 为开始第 $n$ 趟迭代所需的完整数值状态。其物理投影包含 80 层生产网格上的六条廓线：
\begin{equation}
\mathbf{p}^{(n)}
=\left\{m_j,T_j,P_{{\rm gas},j},n_{e,j},
\kappa_{{\rm R},j},g_{{\rm rad},j}\right\}_{j=1}^{80}
\end{equation}
\noindent 而完整求解器状态还携带上一次温度修正、累积的罗斯兰不透明度信息与迭代控制量。记 $\mathcal{G}$ 为作用于该增广状态的一个完整 Kurucz 循环：
\begin{equation}
\mathbf{x}^{(n+1)}=\mathcal{G}\left(\mathbf{x}^{(n)}\right).
\end{equation}
\noindent 它返回更新后的状态，其物理投影是同一恒星标签下重映射后的六场大气。收敛的求解器状态即其不动点：
\begin{equation}
\mathbf{x}_\star=\mathcal{G}\left(\mathbf{x}_\star\right).
\end{equation}
\noindent 初始化就是选择 $\mathbf{x}^{(0)}$。在不动点附近，误差 $\mathbf{e}^{(n)}=\mathbf{x}^{(n)}-\mathbf{x}_\star$ 按下式演化：
\begin{equation}
\mathbf{e}^{(n+1)}
\simeq
\mathbf{J}_{\mathcal{G}}(\mathbf{x}_\star)\mathbf{e}^{(n)} .
\end{equation}
\noindent 局部收敛要求 $\mathbf{J}_{\mathcal{G}}$ 的谱半径小于 1，即一趟遍历缩小而非放大小误差。这个条件定义了一个局部收敛域（basin）。它的边界可能很复杂，因为温度与压强改变布居和不透明度，不透明度改变下一步修正，重映射把光学深度与柱质量耦合，对流边界还可能移动。

\begin{EN}
For fixed requested values of $T_{\rm eff}$, $\log g$, $\xi$, and the
abundances $\{[\mathrm{X}/\mathrm{H}]\}$, let $\mathbf{x}^{(n)}$ denote the
complete numerical state needed to begin iteration $n$. Its physical projection
contains the six profiles on the 80-layer production grid,
\begin{equation}
\mathbf{p}^{(n)}
=\left\{m_j,T_j,P_{{\rm gas},j},n_{e,j},
\kappa_{{\rm R},j},g_{{\rm rad},j}\right\}_{j=1}^{80}
\end{equation}
\noindent while the complete solver state also carries the previous temperature
correction, accumulated Rosseland-opacity information, and iteration controls. We use
$\mathcal{G}$ for one complete Kurucz cycle acting on this augmented state,
\begin{equation}
\mathbf{x}^{(n+1)}=\mathcal{G}\left(\mathbf{x}^{(n)}\right).
\end{equation}
\noindent It returns an updated state whose physical projection is the remapped
six-field atmosphere at the same stellar labels. The converged solver state is
its fixed point,
\begin{equation}
\mathbf{x}_\star=\mathcal{G}\left(\mathbf{x}_\star\right).
\end{equation}
\noindent Initialization means choosing $\mathbf{x}^{(0)}$. Near the fixed point, the error
$\mathbf{e}^{(n)}=\mathbf{x}^{(n)}-\mathbf{x}_\star$ evolves as
\begin{equation}
\mathbf{e}^{(n+1)}
\simeq
\mathbf{J}_{\mathcal{G}}(\mathbf{x}_\star)\mathbf{e}^{(n)} .
\end{equation}
\noindent Local convergence requires the spectral radius of
$\mathbf{J}_{\mathcal{G}}$ to be below unity, so one pass reduces rather than
amplifies a small error. This condition defines a local convergence basin.
Its boundary can be complicated because temperature and pressure change the
populations and opacity, opacity changes the next correction, remapping couples
optical depth to column mass, and the convective boundary can move.
\end{EN}

不动点残差
\begin{equation}
\mathcal{R}(\mathbf{x})
=\mathcal{G}(\mathbf{x})-\mathbf{x}
\end{equation}
\noindent 衡量完整携带状态在一趟遍历下是否自洽。起始与收敛的六条廓线很接近，并不保证完整状态在下一趟下变化很小：小的廓线误差仍可能叠加成大的压强、不透明度或对流响应。因此一个有用的初始化器必须落在物理迭代的收敛域之内，而不仅仅是最小化廓线误差。

\begin{EN}
The fixed point residual
\begin{equation}
\mathcal{R}(\mathbf{x})
=\mathcal{G}(\mathbf{x})-\mathbf{x}
\end{equation}
\noindent measures whether the complete carried state is consistent under one
pass. Close agreement between the six starting and converged profiles does not
guarantee that the complete state changes little under the next pass. Small
profile errors can still combine into a large pressure, opacity, or convection
response. A useful initializer must therefore lie inside the basin of the
physical iteration, not merely minimize profile error.
\end{EN}

\begin{physbox}{不动点、雅可比谱半径与收敛域}
把大气迭代写成 $\mathbf{x}^{(n+1)}=\mathcal{G}(\mathbf{x}^{(n)})$ 后，收敛问题就变成标准的\textbf{不动点迭代}理论：解 $\mathbf{x}_\star$ 满足 $\mathcal{G}(\mathbf{x}_\star)=\mathbf{x}_\star$；在解附近线性化，误差每趟乘以雅可比矩阵 $\mathbf{J}_{\mathcal{G}}$。若其\textbf{谱半径}（最大特征值的模）小于 1，误差逐趟衰减——初值就在收敛域内；若大于 1，再接近的初值也可能被弹开。这给“什么样的初值算好”一个精确判据：不是和答案长得像，而是落在收敛域里。后文讨论冷星发散时（第 8 节），物理图像就是小温度扰动会重组分子平衡与对流边界，等效于谱半径超过 1，固定点迭代失去局部稳定性。
\end{physbox}

爱丁顿灰色大气提供了一种通用的初始结构，但它缺乏对谱线覆盖、分子平衡与对流的响应 \citep{Mihalas1978,Gustafsson2008}，因此可能落在有用的收敛域之外。预计算大气网格能提供邻近的收敛起始模型 \citep{CastelliKurucz2003}，但存档模型只有在物理迭代仍处于目标解收敛域内时才有用。

\begin{EN}
An Eddington grey atmosphere provides a generic starting structure, but it
lacks the response to line blanketing, molecular equilibrium, and convection
\citep{Mihalas1978,Gustafsson2008}. It can therefore lie outside the useful
basin. A precomputed atmosphere grid supplies nearby converged starting models
\citep{CastelliKurucz2003}, although a stored model is useful only when the
physical iteration remains in the target solution's basin.
\end{EN}

\begin{physbox}{灰色大气是什么}
“灰色大气”是最简单的解析大气模型：假设不透明度与波长无关（“灰色”），且能量全靠辐射扩散输运，可解出 $T(\tau_{\rm R})=T_{\rm eff}[\frac{3}{4}(\tau_{\rm R}+\frac{2}{3})]^{1/4}$。它抓住了“越深越热”的整体趋势，但真实大气的不透明度强烈依赖波长（谱线覆盖），冷星还有分子与对流。灰色初值与真实结构差得远，物理迭代要跑很多趟甚至失败——这就是需要一个“懂物理的”学习型初值的原因。后文把温度归一化到灰色廓线上再做预测，正是为了让网络只学习相对灰色背景的\textbf{偏离}部分，降低学习难度。
\end{physbox}

训练初始化器首先需要一批收敛的物理大气。更快的物理迭代使这批初始集合可以负担。之后交替进行物理求解与重训练来扩充语料，只有收敛的解才被保留为训练目标。

\begin{EN}
Training an initializer first requires a set of converged physical atmospheres.
Faster physical iteration makes this initial set affordable. Alternating
physical solutions and retraining then expands the corpus, with only converged
solutions retained as training targets.
\end{EN}

\subsection{紧凑的物理大气表示 A compact physical atmosphere representation}

核心设计选择是同时预测全部六个场。它们的原始数值跨越多个数量级，且满足正性与单调性约束。我们把第 $j$ 层变换为
\begin{equation}
\begin{aligned}
\mathbf{u}_j=\bigg(&
\log_{10}\Delta m_j,\,
\log_{10}\frac{T_j}{T_{{\rm grey},j}},\,
\log_{10}P_{{\rm gas},j},\\
&\log_{10}n_{e,j},\,
\log_{10}\kappa_{{\rm R},j},\,
\operatorname{asinh}\frac{g_{{\rm rad},j}}{s_g}
\bigg),
\end{aligned}
\end{equation}
\noindent 其中 $\Delta m_1=m_1$、$j>1$ 时 $\Delta m_j=m_j-m_{j-1}$，且
\begin{equation}
T_{\rm grey}(\tau_{\rm R})=
T_{\rm eff}
\left[\frac{3}{4}\left(\tau_{\rm R}+\frac{2}{3}\right)\right]^{1/4}.
\end{equation}
\noindent 正的质量增量保证单调性，对数保证正性，灰色归一化去掉了大部分整体温度趋势。家族专属的标度 $s_g$ 是训练语料中辐射加速度绝对值的中位数；带符号变换保留了辐射加速度的两个方向。

\begin{EN}
The central design choice is to predict all six fields together. Their raw
values span many orders of magnitude and obey positivity and monotonicity
constraints. We therefore transform layer $j$ to
\begin{equation}
\begin{aligned}
\mathbf{u}_j=\bigg(&
\log_{10}\Delta m_j,\,
\log_{10}\frac{T_j}{T_{{\rm grey},j}},\,
\log_{10}P_{{\rm gas},j},\\
&\log_{10}n_{e,j},\,
\log_{10}\kappa_{{\rm R},j},\,
\operatorname{asinh}\frac{g_{{\rm rad},j}}{s_g}
\bigg),
\end{aligned}
\end{equation}
\noindent where $\Delta m_1=m_1$ and
$\Delta m_j=m_j-m_{j-1}$ for $j>1$, and
\begin{equation}
T_{\rm grey}(\tau_{\rm R})=
T_{\rm eff}
\left[\frac{3}{4}\left(\tau_{\rm R}+\frac{2}{3}\right)\right]^{1/4}.
\end{equation}
\noindent Positive mass increments enforce monotonicity, logarithms preserve
positivity, and the grey scaling removes most of the global temperature trend.
The family-specific scale $s_g$ is the median absolute radiative acceleration
in the training corpus. The signed transform retains both directions of
radiative acceleration.
\end{EN}

变换后的 $80\times6$ 大气共 480 个坐标。联合主成分分析（PCA）\citep{JolliffeCadima2016} 把它压缩为 $K=160$ 个系数，同时保留超过 99.999\% 的标准化训练方差。

\begin{EN}
The transformed $80\times6$ atmosphere contains 480 coordinates. A joint
principal component analysis, or PCA \citep{JolliffeCadima2016}, reduces it to
$K=160$ coefficients while retaining more than 99.999 percent of the
standardized training variance.
\end{EN}

一个前馈网络把表~\ref{tab:initializer_families} 中的五标签、八标签输入映射到标准化的 PCA 系数。对直接丰度家族，一个元素感知的集合编码器先把每种丰度与其元素身份结合 \citep{Zaheer2017}，再把该摘要与 $T_{\rm eff}$、$\log g$、$\xi$（有效温度、表面重力、微观湍流）拼接，然后预测同样的系数表示。若 $\widehat{\mathbf z}$ 为标准化网络输出，解码则逆转系数与廓线的两步标准化：
\begin{equation}
\begin{aligned}
\widehat{\mathbf c}
&=\overline{\mathbf c}+\mathbf{s}_{c}\odot\widehat{\mathbf z},
\\
\operatorname{vec}(\widehat{\mathbf u})
&=\overline{\mathbf u}
+\mathbf{s}_{u}\odot\mathbf{B}_{K}^{\mathsf T}\widehat{\mathbf c},
\\
\mathbf{p}^{(0)}&=\mathcal{D}(\widehat{\mathbf{u}}).
\end{aligned}
\end{equation}
\noindent 其中 $\overline{\mathbf c}$、$\mathbf{s}_c$ 为 PCA 系数的训练均值与标度；对应的廓线量为 $\overline{\mathbf u}$、$\mathbf{s}_u$；$\mathbf{B}_K$ 为保留的 $K$ 分量 PCA 基；帽子标记网络预测；$\operatorname{vec}$ 把深度×场数组拉平；$\odot$ 为逐元素乘；$\mathsf T$ 为转置。解码器 $\mathcal{D}$ 逆转物理变换，返回六条起始廓线 $\mathbf{p}^{(0)}$，求解器再围绕这些廓线构建其余携带状态。联合基保留了跨场、跨深度的相关变化——这很重要，因为物理求解器对六条廓线是整体响应的。

\begin{EN}
A feedforward network maps the five- and eight-label inputs in
Table~\ref{tab:initializer_families} to standardized PCA coefficients. For the
direct-abundance family, an element-aware set encoder first combines each
abundance with its element identity \citep{Zaheer2017}, then joins that summary to
$T_{\rm eff}$, $\log g$, and $\xi$, which denote effective temperature,
surface gravity, and microturbulence, before
predicting the same coefficient representation. If $\widehat{\mathbf z}$ is
the standardized network output, decoding reverses both coefficient and
profile standardization,
\begin{equation}
\begin{aligned}
\widehat{\mathbf c}
&=\overline{\mathbf c}+\mathbf{s}_{c}\odot\widehat{\mathbf z},
\\
\operatorname{vec}(\widehat{\mathbf u})
&=\overline{\mathbf u}
+\mathbf{s}_{u}\odot\mathbf{B}_{K}^{\mathsf T}\widehat{\mathbf c},
\\
\mathbf{p}^{(0)}&=\mathcal{D}(\widehat{\mathbf{u}}).
\end{aligned}
\end{equation}
\noindent Here $\overline{\mathbf c}$ and $\mathbf{s}_c$ are the training mean
and scale of the PCA coefficients. The corresponding profile quantities are
$\overline{\mathbf u}$ and $\mathbf{s}_u$, while $\mathbf{B}_K$ is the retained
$K$-component PCA basis. Hats identify network predictions.
$\operatorname{vec}$ flattens the depth-by-field profile array, $\odot$
denotes elementwise multiplication, and $\mathsf T$ denotes transpose. The
decoder $\mathcal{D}$ reverses the physical transforms and returns the six
starting profiles $\mathbf{p}^{(0)}$. The solver constructs the remaining
carried state around these profiles.
A joint basis preserves correlated variations across fields and depth, which
matters because the physical solver responds to all six profiles together.
\end{EN}

\begin{physbox}{为什么要做这些变换和 PCA}
\begin{itemize}
\item \textbf{变换的目的}：神经网络最擅长拟合量级相近、平滑变化的量。大气各场跨越十几个数量级（如 $n_e$ 从 $10^8$ 到 $10^{18}$），直接预测会被大数主导、小数失真。取对数把乘法变加法；预测 $\Delta m$ 而非 $m$ 自动保证柱质量单调递增；温度除以灰色廓线则把“必然存在的整体趋势”剥离，网络只学物理上有趣的偏离。
\item \textbf{PCA 的作用}：480 个坐标并非独立——相邻层温度相关，温度与压强被流体静力学绑定。PCA 找到这些相关性的主轴，用 160 个系数即可重建 99.999\% 的方差。网络只需预测 160 个数，训练更稳、更省样本。
\item \textbf{物理约束损失}（下式）：损失函数不只看廓线像不像，还检查深度导数、$\mathrm{d}\tau_{\rm R}/\mathrm{d}m=\kappa_{\rm R}$ 和流体静力学平衡——即预测结构要满足大气成立的基本方程。这让初值更可能落在物理迭代的收敛域内。
\end{itemize}
\end{physbox}

我们在 PCA 解码之后评估训练损失，因此它约束的是完整深度结构而非仅系数误差。四项分别对应不同的重启误差：$\mathcal{L}_{\rm prof}$ 匹配六条解码廓线的数值；$\mathcal{L}_{\nabla}$ 匹配相邻层之间的变化，防止整体水平正确但深度依赖错误；$\mathcal{L}_{\tau}$ 检查柱质量与罗斯兰光学深度的映射；$\mathcal{L}_{\rm hse}$ 检查流体静力学压强平衡。系数 $\lambda$ 设定后三项的相对权重：
\begin{equation}
\mathcal{L}
=\mathcal{L}_{\rm prof}
+\lambda_{\nabla}\mathcal{L}_{\nabla}
+\lambda_{\tau}\mathcal{L}_{\tau}
+\lambda_{\rm hse}\mathcal{L}_{\rm hse}.
\end{equation}
\noindent 廓线与梯度项对变换后的场使用稳健的 smooth-L1 失配；光学深度与流体静力学项在以物理单位重建温度、压强、不透明度与柱质量之后求值，衡量对下式的偏离：
\begin{equation}
\frac{\mathrm{d}\tau_{\rm R}}{\mathrm{d}m}=\kappa_{\rm R},
\qquad
\frac{\mathrm{d}P_{\rm gas}}{\mathrm{d}m}\simeq g-g_{\rm rad}.
\end{equation}
\noindent 光学深度项把积分的 $\kappa_{\rm R}\,\mathrm{d}m$ 关系与收敛目标比较；流体静力学项把气压梯度与 $g-g_{\rm rad}$ 比较，并惩罚气压随深度下降的情况。这些选择保留了强烈影响物理重启的廓线形状与两条结构平衡，但不替代辐射与对流平衡——那仍须由物理求解器建立。五标签与八标签模型取 $\lambda_{\nabla}=0.1$、$\lambda_{\tau}=\lambda_{\rm hse}=0.05$；直接丰度模型取 $\lambda_{\nabla}=0.2$，物理一致性项权重相同。

\begin{EN}
We evaluate the training loss after PCA decoding, so it constrains the full
depth structure rather than coefficient errors alone. Four terms address
distinct restart errors. $\mathcal{L}_{\rm prof}$ matches the values of all six
decoded profiles. $\mathcal{L}_{\nabla}$ matches their changes between adjacent
layers, which prevents a profile with the correct overall level from retaining
the wrong depth dependence. $\mathcal{L}_{\tau}$ checks the mapping between
column mass and Rosseland optical depth, while $\mathcal{L}_{\rm hse}$ checks
hydrostatic pressure balance. The coefficients $\lambda$ set the relative
weight of the last three terms,
\begin{equation}
\label{eq:initializer_loss}
\mathcal{L}
=\mathcal{L}_{\rm prof}
+\lambda_{\nabla}\mathcal{L}_{\nabla}
+\lambda_{\tau}\mathcal{L}_{\tau}
+\lambda_{\rm hse}\mathcal{L}_{\rm hse}.
\end{equation}
\noindent The profile and gradient terms use a robust smooth-L1 mismatch on
the transformed fields. The optical-depth and hydrostatic terms are evaluated
after reconstructing temperature, pressure, opacity, and column mass in
physical units. They measure departures from
\begin{equation}
\frac{\mathrm{d}\tau_{\rm R}}{\mathrm{d}m}=\kappa_{\rm R},
\qquad
\frac{\mathrm{d}P_{\rm gas}}{\mathrm{d}m}\simeq g-g_{\rm rad}.
\end{equation}
\noindent The optical-depth term compares the integrated
$\kappa_{\rm R}\,\mathrm{d}m$ relation with the converged target. The
hydrostatic term compares the gas-pressure gradient with $g-g_{\rm rad}$ and
also penalizes decreasing gas pressure. These choices preserve the profile
shape and two structural balances that strongly affect a physical restart.
They do not replace radiative and convective equilibrium, which the physical
solver must still establish. The five- and eight-label
models use $\lambda_{\nabla}=0.1$ and
$\lambda_{\tau}=\lambda_{\rm hse}=0.05$. The direct-abundance model uses
$\lambda_{\nabla}=0.2$ with the same weights on the physical-consistency
terms.
\end{EN}

\begin{figure}[t]
\centering
\includegraphics[width=0.98\textwidth]{figure_04.pdf}
\caption{物理大气结构与学习型初始化。各列分别改变 $T_{\rm eff}$、$\log g$、$[\mathrm{M}/\mathrm{H}]$、$[\alpha/\mathrm{M}]$ 或微观湍流（按色标）；各行为公共罗斯兰光学深度网格上的六条变换后深度廓线。虚线为收敛的物理大气，实线为用于初始化物理求解器的八标签预测。（Physical atmosphere structure and learned initialization.）}
\label{fig:initializer_profiles}
\end{figure}

\subsection{自举、标签覆盖与验证 Bootstrap, label coverage, and validation}\label{sec:atmosphere_emulator_bootstrap}

我们用 Kurucz-A1 初始化器 \citep{Li2025} 的起始猜测播下第一个标签块的种子。只有经物理求解器收敛的大气才进入语料；这些解训练出一个临时初始化器，其预测减少了下一块所需的物理趟数；如此反复求解与重训练，得到表~\ref{tab:initializer_families} 中的语料。预测值从未被用作训练目标。

\begin{EN}
We seeded a first label block with starting guesses from the Kurucz-A1
initializer \citep{Li2025}. Only atmospheres that converged with the physical
solver entered the corpus. Those solutions trained a provisional initializer,
whose predictions reduced the physical passes needed for the next block.
Repeating physical solves and retraining produced the corpora summarized in
Table~\ref{tab:initializer_families}. A prediction never served as a training
target.
\end{EN}

发布的三个家族是实用的默认选择，而非对方法的限制。每个训练大气只需数秒到数十秒，巡天项目完全可以针对自己的标签范围或丰度参数化并行生成数万个。

\begin{EN}
The three released families are practical defaults rather than restrictions on
the method. With each training atmosphere taking seconds to tens of seconds, a
survey can generate tens of thousands in parallel for its own label range or
abundance parameterization.
\end{EN}

五标签家族覆盖以整体金属丰度和 $\alpha$ 模式足以描述大气的广泛恒星工作。拉丁超立方边际抽样 \citep{McKay1979} 在不使用指数级张量网格的情况下对每个标签范围分层。对 $T_{\rm eff}$ 与 $\log g$ 秩次施加的弱相关性让常见的矮星与巨星结构被更密地采样，而不把语料限制在某条演化轨迹上。

\begin{EN}
The five-label family covers broad stellar work in which a bulk metallicity and an
alpha pattern describe the atmosphere adequately. Latin-hypercube
marginals \citep{McKay1979} stratify every label range without an exponentially
large tensor grid. A small imposed correlation between the $T_{\rm eff}$ and
$\log g$ ranks samples common dwarf and giant structures more densely without
restricting the corpus to an evolutionary track.
\end{EN}

八标签家族面向演化恒星的丰度应用：挖掘上翻（dredge-up）使表面 C、N 独立于整体金属与 $\alpha$ 模式而变化 \citep{Salaris2015,Martig2016}。这些变化重新分配 CO、CN、OH 中的原子，可改变分子不透明度与大气结构 \citep{Tsuji1973,Gustafsson2008}。C、N、O 的条件拉丁超立方面板附着在抽样得到的五标签母状态上，每个八标签大气独立求解。在两个低维家族中，$[\alpha/\mathrm{M}]$ 在整体 $[\mathrm{M}/\mathrm{H}]$ 缩放之上叠加于 $\alpha$ 元素；CNO 家族中显式的 $[\mathrm{O}/\mathrm{M}]$ 坐标取代氧的这一缩放，而 $[\alpha/\mathrm{M}]$ 继续设定 Ne、Mg、Si、S、Ca、Ti。

\begin{EN}
The eight-label family addresses abundance applications in evolved stars,
where dredge-up changes surface C and N independently of the bulk metal and
alpha patterns \citep{Salaris2015,Martig2016}. These changes redistribute
atoms among CO, CN, and OH and can alter molecular opacity and atmospheric
structure \citep{Tsuji1973,Gustafsson2008}. Conditional Latin-hypercube
panels in C, N, and O are attached to sampled five-label parent states, and
every resulting eight-label atmosphere is solved independently.
In both lower-dimensional families, $[\alpha/\mathrm{M}]$ adds to the bulk
$[\mathrm{M}/\mathrm{H}]$ scaling for the alpha elements. The explicit
$[\mathrm{O}/\mathrm{M}]$ coordinate replaces that scaling for oxygen in the
CNO family, while $[\alpha/\mathrm{M}]$ continues to set Ne, Mg, Si, S, Ca,
and Ti.
\end{EN}

\begin{physbox}{为什么巨星需要单独的 C、N、O 坐标}
恒星离开主序后会经历对流“挖掘上翻”：对流包层加深，把核心氢燃烧（CNO 循环）的加工产物翻到表面，使 C 减少、N 增多——这颗星的表面 C、N 不再反映其出生时的整体金属丰度，而是携带其质量与演化阶段的信息（C/N 比甚至被用作红巨星的“时钟”）。因此对演化恒星，一个整体 $[\mathrm{M}/\mathrm{H}]$ 加 $[\alpha/\mathrm{M}]$ 的两参数模式无法描述真实大气；C、N、O 必须成为独立坐标。又因为 CO 分子的结合能极高，C 与 O 谁少谁就被锁进 CO，剩下的参与 CN、OH 等分子——所以 C/O 比的细微变化会显著改变分子不透明度，这也是大气模型必须自洽处理分子平衡的原因。
\end{physbox}

直接丰度家族接受 81 个独立变化的 $[\mathrm{X}/\mathrm{H}]$ 丰度，覆盖 Li 到 Bi（除 Tc、Pm）以及 Th、U，没有整体金属丰度或 $\alpha$ 坐标。网络把同一请求混合物重述为 $[\mathrm{Fe}/\mathrm{H}]$ 与 80 个非铁 $[\mathrm{X}/\mathrm{Fe}]$ 值，但这一坐标变换不改变传给物理求解器的丰度。

\begin{EN}
The direct-abundance family accepts 81 independently variable
$[\mathrm{X}/\mathrm{H}]$ abundances. It covers Li through Bi except Tc and Pm,
together with Th and U. It has no bulk
metallicity or alpha coordinate. The network re-expresses the same requested
mixture as $[\mathrm{Fe}/\mathrm{H}]$ and 80 non-iron
$[\mathrm{X}/\mathrm{Fe}]$ values, but this change of coordinates does not alter
the abundances passed to the physical solver.
\end{EN}

我们首先在六个光谱对照上验证五标签与八标签家族：热矮星、太阳、红巨星，以及富碳、富金属、富氧混合物。单元素覆盖值是绝对的 $[\mathrm{X}/\mathrm{H}]$，替换该元素的整体缩放。

\begin{EN}
We first validate the five- and eight-label families on six spectral controls spanning a hot
dwarf, the Sun, a red giant, and C-enhanced, metal-rich, and O-enhanced
mixtures. Individual abundance
overrides are absolute $[\mathrm{X}/\mathrm{H}]$ values and replace the bulk
scaling for that element.
\end{EN}

图~\ref{fig:initializer_convergence} 左图分离了两种误差。空心符号是 $K=160$ PCA 压缩下限，其标准化坐标 RMS 中位数五标签为 $6.3\times10^{-4}$、八标签为 $4.9\times10^{-4}$——即把 480 个廓线坐标压到 160 个系数再解码，丢弃的大气结构很少。实心符号包含从恒星标签经网络与 PCA 解码的完整预测，对应误差为 $1.5\times10^{-2}$ 与 $1.1\times10^{-2}$，比压缩下限大一个数量级以上。因此当前精度由学习到的“标签到大气”关系决定，而非 PCA 压缩。

\begin{EN}
The left panel of Figure~\ref{fig:initializer_convergence} separates two errors.
Open symbols show the $K=160$ PCA compression floor, whose median
standardized-coordinate RMS is $6.3\times10^{-4}$ for the five-label family and
$4.9\times10^{-4}$ for the eight-label family. Thus compressing the 480 profile
coordinates to 160 coefficients and decoding them discards little of the
atmosphere structure. Filled symbols include the complete prediction from
stellar labels through the network and PCA decoder. Their corresponding errors
are $1.5\times10^{-2}$ and $1.1\times10^{-2}$, more than an order of magnitude
larger than the compression floor. The learned label-to-atmosphere relation,
rather than PCA compression, therefore sets the current accuracy.
\end{EN}

中图问：从每个五/八标签预测出发做三趟求解后，还可能残留多少光谱误差？记 $f_3$、$f_5$、$f_{\rm ref}$ 分别为三趟、五趟后与存档参考的光谱，则 $|f_3-f_{\rm ref}|\leq |f_3-f_5|+|f_5-f_{\rm ref}|$。我们报告这一保守上限的最大值：全部六个对照对两个发布的初始化器都低于 $5.1\times10^{-3}$。三趟只是进度诊断，不是停止规则。

\begin{EN}
The middle panel asks how much spectral error can remain after starting from
each five- and eight-label prediction and taking three solver passes. If
$f_3$, $f_5$, and $f_{\rm ref}$ are the spectra after three passes, after five
passes, and from the stored reference, then
$|f_3-f_{\rm ref}|\leq |f_3-f_5|+|f_5-f_{\rm ref}|$. We report the maximum of
this conservative bound. All six controls lie below $5.1\times10^{-3}$ for
both released initializers. Three passes are a progress diagnostic rather than
a stopping rule.
\end{EN}

右图在 16 个匹配的留出案例中单独考察初始大气的作用。为模拟从预计算网格的邻近节点出发，我们按预先声明的偏移（$T_{\rm eff}$ $\pm250$ K、$\log g$ $\pm0.25$ dex、金属丰度 $\pm0.20$ dex）选取附近的收敛结构，仅用存档结构初始化，然后每颗星都在精确目标标签与丰度下求解。

\begin{EN}
The right panel isolates the effect of the starting atmosphere in 16 matched
held-out cases. To emulate starting from a nearby node in a precomputed grid,
we selected converged structures near two sets of predeclared offsets. They are
$\pm250$ K in $T_{\rm eff}$, $\pm0.25$ dex in $\log g$, and $\pm0.20$ dex in
metallicity. Only the stored structure
initializes the calculation. Every atmosphere is then solved at the exact
target labels and abundances.
\end{EN}

两类网格式起点的迭代中位数约 12 趟，而五标签初始化器为 5.5 趟、八标签为 4 趟——学习型起点把迭代次数分别压缩约 2 倍和 3 倍。结合实测的 6--7 倍单趟加速，八标签计算在这些匹配案例中总共和快约 20 倍。所保留的原程序基准每批请求 30 趟迭代；4 趟物理遍历加更快的单趟，相对该基准约有 50 倍增益，而邻近网格大气给出的是上文中更小、更现实的匹配比较。

\begin{EN}
The median of the two grid-like starts is about 12 iterations, compared with
5.5 for the five-label initializer and 4 for the eight-label initializer. The
learned starts therefore reduce the iteration count by factors of about two and
three. Combined with the measured six- to seven-fold per-iteration gain, the
eight-label calculation is about 20 times faster in these matched cases. The
retained original-program benchmark requests 30 iterations per batch. Four
physical passes and the faster iteration give a gain of about 50 relative to
that benchmark, while a nearby grid atmosphere gives the smaller and more realistic
matched comparison above.
\end{EN}

图~\ref{fig:initializer_profiles} 在廓线空间直接检查八标签初始化器：实线预测与虚线收敛物理廓线对比了 25 个大气（其 C、N、O 随整体混合物变化）。每列改变一个恒星标签，六行展示耦合的温度、压强、密度、不透明度与深度响应。该比较检验从标签经网络与 PCA 解码的完整映射。把标准化残差对所有廓线、大气与深度点合并，全深度 RMS 为 0.028，$\log\tau_{\rm R}\geq1$ 处为 0.027。

\begin{EN}
Figure~\ref{fig:initializer_profiles} provides a direct profile-space check of
the eight-label initializer. Solid predictions are compared with dashed
converged physical profiles for 25 atmospheres whose C, N, and O values follow
the corresponding bulk mixture. Each column varies one stellar label while the
six rows trace the coupled temperature, pressure, density, opacity, and depth
responses. This comparison tests the complete mapping from labels through the
network and PCA decoder. Pooling the standardized residuals over all profiles,
atmospheres, and depth points gives an RMS of 0.028 over the full depth range
and 0.027 for $\log\tau_{\rm R}\geq1$.
\end{EN}

对直接丰度家族，623 个测试对象含 313 个内部与 310 个外部大气。内部大气未参与网络拟合，其温度与气压误差的 95 百分位为 0.67\% 与 0.034 dex。另做的 16 案例重启诊断在 8--14 趟物理遍历内收敛，中位数约 9 趟。在下文应用中，该网络提供初始状态，每个被采用的大气随后都经物理求解与验证。

\begin{EN}
For the direct-abundance family, the 623-object test split contains 313
internal and 310 external atmospheres. The internal atmospheres were excluded
from network fitting and have 95th-percentile temperature and gas-pressure
errors of 0.67 percent and 0.034 dex. A separate 16-case restart diagnostic
converges in 8--14 physical passes, with a median near 9. In the application
below, this network supplies the starting state and every retained atmosphere
is then solved and validated physically.
\end{EN}

这个家族的插值问题更难：3 个恒星参数加 81 个独立丰度构成 84 维输入，训练语料对该空间的采样比五/八标签语料稀疏得多。按化学或不透明度响应对丰度分组的坐标，有望在未来初始化器中降低有效维度。

\begin{EN}
This family presents a harder interpolation problem because three stellar
parameters and 81 independently varying abundances define an 84-dimensional
input. Its training corpus samples that space much more sparsely than the
five- and eight-label corpora. Coordinates that group abundances by their
chemical or opacity response may reduce the effective dimension in future
initializers.
\end{EN}
